\section{讨论与局限}
\label{sec:limitations}
\textbf{路径增量与机制解释。}
匹配路径对照检验的是当前实现、拟合数据和优化预算下，
显式本流读出与时空上下文组合是否值得使用。
Direct访问的真实观测本来就在SPIN可用的输入集合内，
所以其增量不能解释为引入了此前不可获得的信息。
三统计量读出提供一种局部插值偏置，但本轮没有分离这一偏置、有效容量和优化行为的作用。
固定参数初始化、数据与日程提高了比较的可解释性，
仍不能证明另一种经过合理设计的浅层SPIN解码器无法获得相同收益。

\textbf{精度与成本的适用范围。}
相对Context-only增加Direct的成本，与在Direct-only上加入整个上下文网络的成本，
是两个不同的工程问题。
后者即使降低NMAE，也不意味着应在所有吞吐或延迟约束下选用组合。
不同误差指标和掩码子组还可能存在取舍；NRMSE的变化不能未经分解就解释为峰值恢复能力。
本轮没有修改融合权重、增加门控或扩展损失来追求有利排序。

\textbf{评价与优化边界。}
开发段同时承担检查点选择和研究决策，两个WAN的历史后续范围也已被项目访问。
因此，配对方向、末段曲线和大量缺失项均不恢复独立测试身份。
两个种子的样本标准差只描述有限重复；训练轮次彼此相关，
末20轮诊断不构成显著性检验。
统一的40轮低学习率阶段只是一项预先固定的优化诊断，
达到160轮不能证明各架构充分收敛；历史Full未参加该诊断，
其结果不支持对同步记忆在其他配方下的普遍否定。

\textbf{任务与复现边界。}
本研究依赖完整历史监督、固定50槽窗口和20\%当前输入预算，
每流始终有至少四个观测，未评价整流完全失观测或在线因果重建。
GEANT时间和流轴到原始采集记录的映射尚需进一步核实，
不能将槽位距离直接换算成某种网络时间尺度。
新实验保留源文件、数据、初始状态、掩码与检查点校验值，
但训练入口仍依赖本机历史注册资产；公开源码本身不构成独立复现完成。

\textbf{投稿前仍需完成的确认。}
下一步应固定候选及检查点、落实可核验的评价来源，
再与正常工作的SPIN、Linear、旧Direct+Sync及一个相关的SRMF加邻近观测补正参考作有限比较。
Context-only不能替代其中的外部SPIN基线。
若只能复用历史流量，应明确称为冻结后复核；
新掩码检验缺失模式变化，不是新的流量样本。
这些工作用于确认实际方法价值，不需要再次扩张融合模块或消融矩阵。

\section{结论}
本文在浅层SPIN骨架中加入本流双向真实观测的显式统计量读出，
并以独立训练的两个单路径对照检验组合的使用价值。
两套WAN、两个初始化及固定日程下，组合在开发NMAE上具有方向一致的增量，
相对Context-only的批量8推理附加成本较小，末段轨迹也支持排序。
另一方面，组合明显慢于Direct-only，部分辅助误差和单侧支持分组仍存在取舍。
本轮据此支持保留SPIN+Direct作为后续确认的主候选，
并将Full限定为历史误差与成本参考。
该结论补强了组件证据，尚不替代公平外部方法比较、评价来源核验与独立可执行的复现路径。
